%%%%%%%%%%%%%%%%%%%%%%%%%%%%%%%%%%%%%%%%%%%%%%%%%%%%%%%%%%%%%%%%%%%%%%
% Lớp 12 | Ứng dụng đạo hàm để khảo sát hàm số | Giá trị lớn nhất và giá trị nhỏ nhất của hàm số
% Dạng 2D1?3-4: Ứng dụng GTNN, GTLN trong bài toán phương trình, bất phương trình
% Số câu: 1  (sinh tự động bởi tach_cau.py; từ nay thầy sửa trực tiếp file này)
%%%%%%%%%%%%%%%%%%%%%%%%%%%%%%%%%%%%%%%%%%%%%%%%%%%%%%%%%%%%%%%%%%%%%%

%%%%%%%%%%%%%%%% TRẢ LỜI NGẮN %%%%%%%%%%%%%%%%

% nguồn: Phan2_Lop_12_Dai_so_Chuong_1
\begin{bt}%[2D1V3-4]%[ma:9487129c]
	Cho bất phương trình $x^2-(m+1)x-m+2 \geq 0$, với $m$ là tham số thực.  Tìm giá trị lớn nhất của $m$ sao cho bất phương trình đã cho nghiệm đúng với mọi $x \in[0; 4]$.
	\shortans{1} 
	\loigiai{ 
		Với $x \in[0; 4]$, ta  có
		\begin{align*}
			 x^2-(m+1)x-m+2 \geq 0 & \Leftrightarrow x^2-x+2 \geq m(x+1) \\
			& \Leftrightarrow \dfrac{x^2-x+2}{x+1} \geq m \, (\text {do } x+1>0 \forall x \in[0; 4]).\quad  (1)
		\end{align*}
		Xét hàm số $f(x)=\dfrac{x^2-x+2}{x+1}$, $x \neq-1$. Hàm số này  liên tục trên $[0; 4]$.\\
		Ta có  $f'(x)=\dfrac{x^2+2 x-3}{(x+1)^2}$ suy ra $f'(x)=0 \Leftrightarrow x=1$ hoặc $x=-3$ (loại vì $-3 \notin[0; 4]$).\\
		Ta có $f(0)=2$; $f(1)=1$; $f(4)=\dfrac{14}{5}$.\\ 
		Suy ra $\min  \limits_{[0; 4]} f(x)=f(1)=1$.\\
		Do đó bất phương trình (1) nghiệm đúng với mọi $\begin{aligned}
		 x \in[0; 4] &\Leftrightarrow m \leq \min \limits_{[0; 4]} f(x) \leq 1\\
		&\Leftrightarrow m \leq 1. 
	\end{aligned}$\\
	Vậy $m=1$ là giá  trị cần tìm.
 }
\end{bt}
